\section{Rule of Thumb Derivation}
\label{sec:rule}

\subsection{The 20-Units-per-District Heuristic}
\label{sec:rule:heuristic}

The resolution selection rule is:
\[
  \text{Use block-group resolution if and only if } \frac{k}{n_{\rm tracts}} > 0.05.
\]
Equivalently: use tract resolution only if each district contains at least
$n_{\rm tracts}/k \geq 20$ tracts on average.

The threshold of 20 units per district is motivated by three considerations.

\paragraph{Population balance feasibility.}
For a census tract with median population 4,084 and a district with ideal
population $T^* = T/k$, the finest achievable population deviation using
whole-tract assignment is approximately:
\[
  \delta_{\min} \approx \frac{\text{median tract population}}{T^*}
  = \frac{4{,}084}{T/k}.
\]
For $\delta_{\min} \leq 0.5\%$, we need:
\[
  \frac{4{,}084}{T/k} \leq 0.005 \implies T/k \geq 816{,}800.
\]
This would restrict tract-resolution redistricting to states with ideal district
populations above $\sim\!800,000$---essentially only states with 1--2 congressional
seats. For state legislative redistricting, tract resolution is inadequate
for \emph{any} state using the balance criterion alone.

However, this calculation is too pessimistic. In practice, the pipeline
uses a post-processing boundary-swap algorithm that can fine-tune district
boundaries by reassigning individual tracts. The effective minimum deviation
depends on the variance of tract populations, not just the maximum. The 0.5\%
constraint is consistently achievable when each district contains $\geq 20$
tracts, because the central limit theorem ensures that the sum of 20 random
tract populations has small variance relative to its mean.

\paragraph{Compactness optimisation richness.}
The number of distinct district configurations for a district containing
$m$ tracts grows as $O(n^m / k^m)$ under a random graph model. For $m \geq 20$,
the number of configurations is large enough that the METIS edge-cut minimisation
can meaningfully distinguish good from bad boundaries. For $m < 10$,
the algorithm has few configurations to choose among, and the output is
nearly independent of the objective function---compactness is determined
by geography rather than optimisation.

\paragraph{Empirical calibration.}
Companion paper C.1 \citep{dellaLibera2026maup} found that congressional
redistricting at tract resolution (where the typical district contains
$\sim\!100$ tracts) gives PP scores approximately equal to block-group resolution.
Our empirical results (Section~\ref{sec:comparison}) show that PP scores begin
to diverge between tract and block-group resolution when $k/n > 0.05$
($n/k < 20$), with the block-group score being higher. Below $k/n = 0.02$
($n/k > 50$), the two resolutions give essentially identical PP scores.

\subsection{Practical Application}
\label{sec:rule:application}

To apply the rule for a given state and chamber:
\begin{enumerate}
  \item Look up $n_{\rm tracts}$ from the 2020 Census TIGER/Line tract shapefile
    for the state.
  \item Compute $k/n_{\rm tracts}$ for the target chamber.
  \item If $k/n_{\rm tracts} > 0.05$, use \texttt{--resolution block\_group}.
    If $k/n_{\rm tracts} > 0.50$, consider \texttt{--resolution block}.
  \item Otherwise, use \texttt{--resolution tract} (default).
\end{enumerate}

The \texttt{redist} pipeline automates this check when
\texttt{--resolution auto} is specified:
\begin{verbatim}
redist state --state WA --chamber house --resolution auto --year 2020
\end{verbatim}
The auto-selection reads the state's tract count from the configuration file
and applies the $k/n > 0.05$ rule.
